% Generated draft: requires mathematical and source review.
\documentclass[11pt]{article}
\usepackage[margin=25mm]{geometry}
\usepackage{fontspec}
\setmainfont{Latin Modern Roman}
\usepackage{amsmath,amssymb,mathtools,unicode-math}
\setmathfont{Latin Modern Math}
\usepackage{enumitem,needspace,xurl,hyperref}
\setlength{\parindent}{0pt}
\setlength{\parskip}{5pt}
\setlength{\emergencystretch}{3em}
\newcommand{\N}{\mathbb{N}}
\newcommand{\Z}{\mathbb{Z}}
\newcommand{\Q}{\mathbb{Q}}
\newcommand{\R}{\mathbb{R}}
\newcommand{\C}{\mathbb{C}}
\newcommand{\F}{\mathbb{F}}
\newcommand{\A}{\mathbb{A}}
\newcommand{\PP}{\mathbb{P}}
\newcommand{\E}{\mathbb{E}}
\newcommand{\eps}{\varepsilon}
\newcommand{\dd}{\,\mathrm{d}}
\begin{document}
\section*{UnsolvedMath ID: 30001942}
\textbf{Source problem code:} catalogue entry
\subsection{Definitions and mathematical statement}
The supplied question concerns the two-dimensional periodic semi-geostrophic system, denoted by \(\mathrm{SG}\), and its dual system, denoted by \(\mathrm{SGD}\). The record does not supply the equations or define their distributional solution classes.

The intended uniqueness question can therefore only be expressed schematically. For \(X\in\{\mathrm{SG},\mathrm{SGD}\}\), let \(\mathcal S_X(D;I)\) denote the intended class of distributional solutions of system \(X\), with prescribed data \(D\), on a common time interval \(I\), subject to the intended two-dimensional periodic spatial conditions. Is it true, separately for each system, that
\[
\forall U,V\in\mathcal S_X(D;I),\qquad U=V,
\]
for every admissible choice of \(D\) and \(I\)? Here equality means equality of the corresponding solution components as distributions on the common space-time domain.

This is a formulation template, not a fully specified PDE statement: the equations, spatial periodicity convention, time interval, prescribed data, admissibility hypotheses, solution components, and precise distributional formulations must be recovered from the cited source. In particular, no regularity, integrability, convexity, or density assumptions are supplied, and none are imposed here. The question asks for uniqueness within each system separately; the record does not state a uniqueness or equivalence assertion between the two systems.
\subsection{Short English statement}
For each of the two-dimensional periodic semi-geostrophic system and its dual, ask whether two distributional solutions with the same prescribed data must coincide. The supplied record omits the equations and admissible solution classes, so a fully precise uniqueness statement requires further source information.
\subsection{Sources}
\begin{itemize}
\item Local UnsolvedMath registry: supplied record with permanent ID \(30001942\), titled \emph{Uniqueness for Periodic Semigeostrophic Systems}, in the Oberwolfach Reports Open Problems collection. Exact source problem code: \texttt{catalogue entry}. The registry records proposed year 2011 and names no individual proposer.
\item Supplied source citation: \emph{Variational Methods for Evolution} (2012), Oberwolfach Reports, pp.\ 3145--3216, DOI \texttt{10.4171/owr/2011/55}. Source URL and supplied literature-source link, labelled ``original OWR report'': \url{https://doi.org/10.4171/owr/2011/55}.
\item The supplied background describes Oberwolfach Reports as documenting mathematical workshops and their open-problem discussions. The registry labels the statement ``exact'' and reports that its mathematical content was checked against the official report, with notation, context, and question form normalized without changing the source-level claim. These are supplied provenance assertions, not independent verification in this draft.
\item Literature assessment dated exactly \texttt{2026-08-21}: the supplied curation labels the problem open, reports that the source explicitly describes uniqueness of distributional solutions as unknown for both \(\mathrm{SG}\) and \(\mathrm{SGD}\), and reports that no complete resolution was found. This is supplied dated curation, not an independently verified assessment of the literature or current mathematical status.
\end{itemize}
\end{document}
